% Frontier Capability Index (FCI) Methodology for KDD Paper
% This section explains how we selected the Pareto frontier models

\subsection{Model Selection and The Pareto Frontier}

To define the candidate pool, we constructed a \textbf{Frontier Capability Index (FCI)} based on a composite of three rigorous benchmarks: \textbf{GPQA} (Graduate-Level Google-Proof Q\&A), \textbf{LiveBench} (Contamination-resistant evaluation), and \textbf{HLE} (Human Level Evaluation). 

We specifically selected these benchmarks to avoid the \textit{saturation effects} observed in older datasets (e.g., MMLU), where the performance gap between efficient and flagship models has narrowed to negligible levels. By focusing on high-difficulty tasks, our Pareto frontier preserves the \textit{capability differentiation} required to evaluate ``Rational Luxury'' routing strategies.

Candidate models were selected strictly based on their position on the Price-vs-FCI curve (Figure~\ref{fig:pareto_frontier}), ensuring that the router's choice set contained only non-dominated options prior to any exposure to the test distribution.

\subsubsection{FCI Construction}

The Frontier Capability Index is computed as the arithmetic mean of normalized benchmark scores:

\begin{equation}
\text{FCI}_{model} = \frac{1}{3} \left( S^{norm}_{HLE} + S^{norm}_{GPQA} + S^{norm}_{LiveBench} \right)
\end{equation}

where each benchmark score is normalized using min-max scaling across all evaluated models:

\begin{equation}
S^{norm} = \frac{S_{model} - S_{min}}{S_{max} - S_{min}}
\end{equation}

Using the full range of models (including weak models) for normalization ensures robust scaling and prevents negative scores. The normalization ranges for our 39 evaluated models are:

\begin{itemize}
\item \textbf{HLE}: $[0.033, 0.372]$
\item \textbf{GPQA}: $[0.200, 0.910]$  
\item \textbf{LiveBench}: $[0.020, 0.920]$
\end{itemize}

\subsubsection{Pareto Frontier Identification}

From 42 candidate models, 39 had complete benchmark data for FCI computation. We identified 4 Pareto-optimal models by iterating through models sorted by cost and selecting those with strictly higher FCI than all cheaper alternatives:

\begin{table}[h]
\centering
\caption{FCI-Based Pareto Frontier Models}
\label{tab:pareto_models}
\begin{tabular}{lcccc}
\toprule
\textbf{Model} & \textbf{Cost} & \textbf{FCI} & \textbf{HLE} & \textbf{GPQA} \\
 & \textbf{(\$/1M)} & & & \\
\midrule
GPT-OSS-120B & 0.06 & 0.740 & 0.185 & 0.78 \\
Gemini-2.5-Pro & 3.44 & 0.766 & 0.211 & 0.84 \\
Claude-Opus-4.5 & 5.00 & 0.876 & 0.284 & 0.87 \\
Gemini-3-Pro & 7.00 & 1.000 & 0.372 & 0.91 \\
\bottomrule
\end{tabular}
\end{table}

\subsubsection{Why Ultra-Cheap Models Were Excluded}

Several models priced below GPT-OSS-120B (\$0.06/1M) were excluded from the Pareto frontier despite lower cost. Analysis reveals these models are \textit{Pareto-dominated}:

\begin{itemize}
\item \textbf{Ministral-3B} (\$0.08/1M, FCI=0.180): Saves \$0.02 but loses 0.56 FCI points—a 28$\times$ worse value proposition
\item \textbf{Gemma-3-4B-IT} (\$0.085/1M, FCI=0.094): Saves \$0.015 but loses 0.65 FCI points—a 430$\times$ worse value
\item \textbf{Ministral-8B} (\$0.10/1M, FCI=0.240): Costs 67\% more than GPT-OSS-120B but delivers 68\% lower FCI
\end{itemize}

GPT-OSS-120B at \$0.06/1M represents an \textit{exceptional value point}: no model offers better quality at lower cost, making it the natural floor of the frontier.

\subsubsection{Frontier Span}

The selected Pareto frontier spans:
\begin{itemize}
\item \textbf{Cost range}: 117$\times$ (\$0.06 to \$7.00 per 1M tokens)
\item \textbf{FCI range}: 0.740 to 1.000 (35\% improvement)
\item \textbf{Quality differentiation}: Sufficient to demonstrate ``Rational Luxury'' routing across difficulty spectrum
\end{itemize}

This wide span enables the router to make economically meaningful trade-offs: simple prompts route to the \$0.06 model, while complex reasoning tasks justify the 117$\times$ premium for marginal 35\% FCI improvement.

